\subsection{Parts X and XI}\label{d98b}

The Details of the chapters on Parts~X and~XI (Sections~\ref{xa:sec}, \ref{xb:sec}, \ref{d11a:sec} and~\ref{d11b}) were written against the note at commit 18a12f9. Where they found a gap, a compressed step or an imprecise wording, the note was corrected, and the Details describe the corrected note. The table lists these places: the Detail, the place in the note, and what changed.

{\small
\begin{longtable}{@{}>{\raggedright\arraybackslash}p{0.17\textwidth}>{\raggedright\arraybackslash}p{0.27\textwidth}>{\raggedright\arraybackslash}p{0.48\textwidth}@{}}
\toprule
Detail & Place in the note & What changed\\
\midrule\endfirsthead
\toprule
Detail & Place in the note & What changed\\
\midrule\endhead
\bottomrule\endfoot
Detail~\ref{xa:d-ext-reading} & Proposition~\Nref{prop:complex-ext}, statement & The statement now names the coercivity constants and the fixed rate of the differences, as in (R2) and (R3) of the Detail.\\
Detail~\ref{xa:d-cmp} & Section~\Nref{sec:complex-references}, paragraph ``Complex references'', first item & At 18a12f9 the two sizes of a complex perturbation were said to be ``comparable up to a constant''. The constant $C_{\rm cmp}=4$ and the threshold $\eps_{\rm cc}$ are now stated, with a sketch of the proof of the Detail.\\
Details~\ref{xa:d-dag}, \ref{xa:d-identity} & Proposition~\Nref{prop:complex-ext}, step~(iii) of the proof & Step~(iii) is now written in the two stages (iii-a) and (iii-b), with the list (b1)--(b4), the stage (iii-a) on the ball of radius $2\delta_*$, and the branch of $\log\det\Delta_k$. At 18a12f9, (b3) and (b4) were not in the list, although the note cited them ``as in step~(iii)''. The note also states that $E=0$ is needed on the ball of radius $2\delta_*$ used in the proof of Proposition~\Nref{prop:dirichlet-hessian}(c).\\
Detail~\ref{xa:d-identity} & Proposition~\Nref{prop:complex-ext}, step~(iii) of the proof & The step now says that Lemma~\Nref{lem:coercive-general}, which gives $\Rea\Delta_k\ge\frac18c_L$ on the ball of complex references, needs $|J_{U_r}|\le\eps$, which holds at every reference where Proposition~\Nref{prop:one-loop-polymer} is applied (Setting of Section~\Nref{sec:step-quantities}), and that its conditions on $\eps$ hold since $\eps\le\eps_{\max}$.\\
Details~\ref{xa:d-step-i}, \ref{xa:d-deltaP} & Proposition~\Nref{prop:complex-ext}, paragraph ``Smallness'' of the proof & At 18a12f9 the condition $C\delta S_{\gamma/4}\le\frac c2$ of Proposition~\Nref{prop:dirichlet-averaged}(e) and the factor $C_{\rm cmp}$ were not in the list of conditions on $\delta_P$. Both now enter through $\delta_*$, in the first item ($\delta_P\le c''_L$ with $C_{\rm cmp}c''_L\le\delta_*$).\\
Details~\ref{xa:d-step-i}, \ref{xa:d-AB-ref}, \ref{xa:d-hessian}, \ref{xa:d-deltaP} & Proposition~\Nref{prop:dirichlet-hessian}(c), proof; Proposition~\Nref{prop:hessian-general}(b), proof & At 18a12f9 the $L^2$ bound of $d^2\mathcal N(0)$ and the change $C_W\delta$ of $W_2^X$ at complex perturbations were compressed (proof of Lemma~\Nref{lem:hessian-gamma}; Proposition~\Nref{prop:complex-ext}(c) and step~(v) of its proof), and the proof of Proposition~\Nref{prop:hessian-general}(b) addressed only $\omega_0$. The proof of Proposition~\Nref{prop:dirichlet-hessian}(c) now gives the step, with $\delta_*$ fixed once (written out in Detail~\ref{d8:A1}), and the proof of Proposition~\Nref{prop:hessian-general}(b) points to it.\\
Detail~\ref{xa:d-majorant} & Proposition~\Nref{prop:complex-ext}, step~(vii) of the proof & At 18a12f9 the factor $1+C_{\rm cx}\delta_c$ was justified by the sentence ``each of them is obtained \dots\ from a bounded number of input bounds by sums, products and inversions'', which does not cover the walk series of Theorem~\Nref{thm:weak-inverse} and the contractions of Theorem~\Nref{thm:weak-bg} and Lemma~\Nref{lem:Ds}. The lemma on monotone majorants, a new argument of the Detail, and the level-by-level definition of $C_{\rm cx}$ are now in step~(vii).\\
Detail~\ref{xa:d-notext} & Proposition~\Nref{prop:complex-ext}, last paragraph of the statement & At 18a12f9 the statement listed the whole of Lemma~\Nref{lem:div-curvature} as not extended, while the proof of Proposition~\Nref{prop:dirichlet-hessian}(c) used its first part. The statement now reads ``the identity of Lemma~\Nref{lem:div-curvature} (its bound on $W_0$ extends)'', and step~(ix) agrees.\\
Detail~\ref{xa:d-deltaP} & Remark~\Nref{rem:domain-constants}, first item & At 18a12f9 the thresholds defining $\eps_{\max}$ were ``computed with the worst case of their constants over complex references of size $\delta_c\le\bar\delta$''. That reading is undefined for $\delta_P\le\delta_c\le\bar\delta$, and a worst case over $\delta_c<\delta_P$ would make $\eps_{\max}$ depend on $\delta_P$. The remark now computes them with every multiplicative constant doubled and with the fixed rates.\\
Detail~\ref{xa:d-AB-bounds} & Lemma~\Nref{lem:AB-bounds}(d) & At 18a12f9 the first display of (d) bounded the sum of the two differences by $C_{\rm loop}\delta_j|A'|$; the proof gives this bound for each term, and the display now bounds each term. No later use depends on the factor $2$.\\
Detail~\ref{xa:d-sq-setting} & Section~\Nref{sec:step-quantities}, Setting & The margin $r\ge3$, needed for $r_E+3\le r$ with $r_E=0$, was not stated at 18a12f9; the Setting now assumes it.\\
Detail~\ref{xa:d-int} & Proposition~\Nref{prop:interaction-pieces}(b), proof & At 18a12f9 the case $k=0$ ($\eta=1$, $A_*=B$) was not treated, and the sentence on ``the first variation'' said why an expansion fails but not what gives the bound. The proof now treats $k=0$ and states the range argument.\\
Detail~\ref{xa:d-gf} & Proposition~\Nref{prop:gf-pieces}, proof & The condition $c_0\le1/(4C_M)$, implicit at 18a12f9 in ``$c_0$ small'' of Lemma~\Nref{lem:exp-gauge}(c), is now stated.\\
Detail~\ref{xa:d-reblock} & Remark~\Nref{rem:reblock-margins} & At 18a12f9 the tree bound read $\sqrt dM+r'$; it now reads $\sqrt d(M+r')$, since $r'$ is an $\infty$-distance.\\
Detail~\ref{xb:d-k0-hessian} & Proposition~\Nref{prop:first-step}(b), proof, item ``Local gauges'' & At 18a12f9 the lower bound on the Hessian in a local gauge was obtained ``by Lemma~\Nref{lem:cov-expansion}'', which concerns only the covariant derivative. The item now uses Proposition~\Nref{prop:action} at $\eta=1$ (the Hessian is the square of the covariant curl up to $O(\eps)$), the argument of the Detail, and the condition $\eps\le c_5$ is added to those of the first step.\\
Details~\ref{xb:d-k0-thresholds}, \ref{xb:d-k0-list} & Proposition~\Nref{prop:first-step}(b), proof, paragraph ``The gradient term''; proof of~(c) & At 18a12f9 the conditions (FS1)--(FS3) of the first step were stated only through $a^{(0)}_{\rm step}$, that is, for the critical references. The choice of $\eps_a$, $\delta_a$ and $c_a$ at $k=0$, which gives them at the non-critical references, and the condition of item~(7) of Detail~\ref{xb:d-k0-list}, are now stated.\\
Detail~\ref{xb:d-k0-walks} & Proposition~\Nref{prop:first-step}, last paragraph of the proof & At 18a12f9 the rate at $k=0$ was justified only by ``with the same choice of the weakening rate''. The paragraph now gives the reason: one walk, with rate at least $\kappa-2\delta\ge\kappa'_w$.\\
Detail~\ref{xb:d-genstep} & Theorem~\Nref{thm:gen-step}(a), proof & Part~(a) of the Detail (the hypotheses of Theorem~\Nref{thm:step} for the density in generation form) was a new argument at 18a12f9; the proof now contains it.\\
Detail~\ref{d11a:E11} & Theorem~\Nref{thm:ol-chart}, proof, ``All $V$''; Theorem~\Nref{thm:balaban}, proof (choice of $a_1$) & A gap of bookkeeping type at 18a12f9: ``All $V$'' uses Proposition~\Nref{prop:AB}(vi) at $U_r=U_0^\circ$, whose hypothesis $a\le c'_3$ was not implied by the choice of $a_1$. The choice of $a_1$ now contains $B_3a_1\le c'_3$, and the proof of Theorem~\Nref{thm:ol-chart} uses it. The main line uses Theorem~\Nref{thm:ol-chart} only near $V=1$, where the condition is not needed.\\
Detail~\ref{d11a:A6} & Proposition~\Nref{prop:ol-laplace}, proof, item~(v) & The continuity argument for $\mathcal P_\delta(V)\subset\mathfrak B(V)\cap\mathfrak U(\theta)$ (Step~7 of the Detail) is now in the note.\\
Detail~\ref{d11a:E13} & Lemma~\Nref{lem:ol-offshell}(d) & The last claim of~(d) is now written in the note.\\
Detail~\ref{d11a:F16} & Remark~\Nref{rem:ol-symmetry} & The transversality used in the remark is now written in the note.\\
Detail~\ref{d11b:beta-identification} & Proposition~\Nref{prop:beta-identification}, proof, the isotropy step & The step now supplies the transversality of $\Pi_\Phi$ from Remark~\Nref{rem:marginal-orders} (Step~5 of the Detail).\\
Detail~\ref{d11b:gl-flatdata} & The paragraph after Definition~\Nref{def:gl-flat-data} & The paragraph now cites the equality $d_k(\pi_n\tilde Y)=d_k(\tilde Y)$ from the proof of Lemma~\Nref{lem:gl-limits}(a).\\
Detail~\ref{d11b:F3} & Proposition~\Nref{prop:flow-F}, (F3) & The constant in $C_3$ is now named $C_{\rm pure}$, a constant with $r_t^{-2}\le C_{\rm pure}p_0^6$.\\
Detail~\ref{d11b:ol-uv} & Theorem~\Nref{thm:ol-uv}, proof & The proof now says that the constant of Theorem~\Nref{thm:ol-cumulative} does not depend on $K$ or on the torus.\\
Detail~\ref{xa:d-step-v} & Proposition~\Nref{prop:complex-ext}, step~(v) of the proof & A correction of this volume (Yoh integrated audit, Round~A, 7 November 2026): the curvature formula lacked the commutator term; it now has the remainder $O(|A_0|^2+\eta|A_0||\partial A_0|)$. The bound $|F_U|\le C(\eps+\delta_c)$ and the note are unchanged.\\
Detail~\ref{xa:d-conv} & Lemma~\Nref{lem:norm-conversion} & A correction of this volume (Round~A): the Detail concluded $C_{\rm conv}=2+C\eps$, but the three suprema of the lemma may be attained on different blocks; it now gives $C_{\rm conv}=4+3C\eps$. The lemma has an unspecified constant and is unchanged.\\
Detail~\ref{xa:d-sq-out} & Remark~\Nref{rem:step-quantities-output}(c) & Round~A: the compact condition $gC_Tp\le c\min(r_I,L^{-d}\zeta)$ follows from (Q1) with $\zeta=w$ only for $c\ge\max(\frac12,C_w^{-1})$, the reverse of what the Detail said. The note and the Detail now state the two conditions $gC_Tp\le\frac12r_I$ and $C_wL^dgC_Tp\le\zeta$.\\
Details~\ref{xb:d-k0-hessian}, \ref{xb:d-k0-lower} & Proposition~\Nref{prop:first-step}(b), proof & A correction of this volume (Round~A): Detail~\ref{xb:d-k0-hessian} squared a reverse triangle inequality without its positive part; the square is now expanded directly, and the Detail is stated for any point of the torus. Detail~\ref{xb:d-k0-lower} said that $M'_0\ge\max(M,L)$ depends only on $d$ and $L$; now $M'_0=\max(M,L,R_{\rm IMS})$. Its step ``Local gauges'' takes the gauge on a cube of radius $40M'+1$ around a lattice point near each centre. The note is unchanged.\\
Details~\ref{d11a:F14}, \ref{d11a:F15} & Lemma~\Nref{lem:ol-remainders}(a), Proposition~\Nref{prop:ol-reduction} & A correction of this volume (Round~A): the rank $\dim\mathfrak g$ per cell holds for each term of the Leibniz expansions of the scalar insertions of kind~(ii); a complete insertion has rank at most $C\dim\mathfrak g$, as the lemma states.\\
Detail~\ref{d11b:ol-linear-average} & Lemma~\Nref{lem:ol-linear-average}, proof & A correction of this volume (Round~A): the gauge transform of the trivial configuration equals $e^{\ii\eta t\,\partial\lambda(b)}$ only up to $O(t^2)$ in the exponent, since the values at the two ends of a bond need not commute; only the first order in $t$ is used.\\
Detail~\ref{d11b:gl-inventory} & Lemma~\Nref{lem:gl-inventory}(ii) & A correction of this volume (Round~A): $n_p$ is $s_p$ minus the contributions to $S(U_0)$, $\ip A{J_{U_0}}$ and $\frac12\|D_{U_0}A\|^2$ of Proposition~\Nref{prop:action}, not minus all its parts of order at most $2$; a quadratic part that depends on the curvature remains.\\
Section~\ref{xa:sec-imp} & none & A correction of this volume (Yoh integrated audit, Round~B, 7 November 2026): the row of Theorem~\Nref{thm:weak-inverse} in the table of imported results now lists the hypothesis of Lemma~\Nref{lem:step-small} and, for~(b), the comparison of the local commutators.\\
Detail~\ref{xb:d-k0-walks} & Proposition~\Nref{prop:first-step}(c), proof & Round~B: the comparison of the local commutators assumed in Theorem~\Nref{thm:weak-inverse}(b) is proved at $k=0$, from the comparison of domains of (W2) and the locality of $\Delta_0$. Claim~(a) of the Detail names it (a statement edit of this volume).\\
Detail~\ref{d11a:E5} & Proposition~\Nref{prop:ol-FP}, proof & A correction of this volume (second integrated audit, Round~D, 8 November 2026, XI-N01): the coarea measure $\sigma_\varphi=\delta(\varphi(V'))\,dV'$ was written without the Haar density $\sigma(0)^{|\mathcal B_\T|}$ of $dV'$; the constant is $c_\T=\sigma(0)^{|\T|}$, not $\sigma(0)^{|\T|-|\mathcal B_\T|}$. Only the $V$-independent constant $c_0$ changes. The note's proof now says at its start that its identities of measures hold up to positive factors independent of $V$, which make up $c_0$.\\
Detail~\ref{d11b:kappa} & Section~\Nref{sec:ol-feynman}, \Neqref{eq:ol-kappa} & A correction of this volume (Round~D, XI-N02): the proof summed over an orthonormal basis of the trace form but used the sums of the generators normalized by $\operatorname{tr}(T^aT^b)=T_\rho\delta^{ab}$; it now names these generators. $\kappa_\rho=C_A/T_\rho$ is unchanged.\\
Detail~\ref{xa:d-ext-reading} & Proposition~\Nref{prop:complex-ext}, statement & The November~10 audit, Round~F (10 November 2026), F-R04-005. The fixed rate of the differences is a quarter of the minimum of the perturbation rate and the two chosen inverse decay rates (Lemma~\Nref{lem:perturbed-inverse}, corrected in both notes). (R2) of the Detail now says that $\gamma_{\rm c}/4$ is this rate, since $\gamma_{\rm c}\le\gamma/2$. No number changed.\\
Detail~\ref{xa:d-AB-bounds} & Lemma~\Nref{lem:AB-bounds}; proof of Proposition~\Nref{prop:AB}(iv) & Round~F, F-R06-005 and F-R06-006. The proof said that $Q'_kQ'^*_k\ne I$ for complex $U$; the identity need not hold, while the $\mathrm b$-transpose identity always holds. The note now says so. A correction of this volume: the remainder of the averages was written $C_i$, a name used for other constants; it now has its own symbol.\\
Section~\ref{xa:sec-imp} & Lemma~\Nref{lem:one-step-expansion} & Round~F, F-R06-004. The row of the lemma in the table of imported results adds $\delta'=\delta_j\le1$, which gives the condition $a\delta'^2\le c_2$ of the corrected paragraph ``Other references''.\\
Details~\ref{xb:d-symflow}, \ref{xb:d-genstep} & Remark~\Nref{rem:symmetrization-flow}(a),(b); proof of Theorem~\Nref{thm:gen-step}(a) & Round~F, F-R09-002. The composed configuration $U_B=e^{\ii\eta A_*(B)}U'$ need not lie in the minimal orbit of $e^{\ii B}W$; its real gauge transform $U_*=U_B^{w_B}$ does, and $w_B$ need not lie in $\mathfrak G_0$. The note and Detail~\ref{xb:d-genstep}(b) now use this fibre statement for $U_*$, transfer the domains $\mathcal U^\flat(Y)$ and the activities to $U_B$ by real gauge invariance, and apply the symmetrization lemma with the real seed $\hat V_B$; no complex minimization is used. Detail~\ref{xb:d-symflow} admits the real gauge transforms of real background fields.\\
Detail~\ref{d11a:D2} & Lemma~\Nref{lem:ol-pullback} & Round~F, F-R09-010. A correction of this volume: the flat reference used $\eps=0$ as a regularity parameter, outside the range $(0,1]$. Now $U_0=1$ is $(M,\eps)$-regular for every $\eps\in(0,1]$, $A_*$ does not depend on $\eps$ by uniqueness, and the bounds, uniform in $\eps$, are used as $\eps\to0$.\\
Detail~\ref{d11b:gl-inventory} & Lemma~\Nref{lem:gl-inventory}(ii) & Round~F, F-R07-006. A correction of this volume: the coordinate formula of the gradient of $\mathcal N$ lacked the factor $\eta^{-d}$ of the pairing of gauge fields.\\
\end{longtable}
}
